\textbf{COLUNI 2023 -- Questão 26}\par\medskip Há muito tempo, a questão em torno de gasodutos na Europa ganha destaques nos noticiários:

\par

Texto originalmente publicado pela Super em 2014: \textbf{Gás natural: o grande motivo por trás do conflito entre Rússia e Ucrânia} Disponível em: https://super.abril.com.br/historia/a-verdade-sobre-a-ucrania/. Acesso em: 11 ago. 2022.

\par

Texto originalmente publicado pela BBC em 2022: \textbf{Rússia x Ucrânia: o gasoduto vital ameaçado pelo conflito} Disponível em: https://www.bbc.com/portuguese/internacional-60317372. Acesso em: 11 ago. 2022.

\par

Tal tema é relevante devido à importância do gás natural como fonte de energia para o funcionamento de indústrias e para o aquecimento das residências. O gás natural é um combustível fóssil, não renovável e, comparadamente, mais limpo do que outros do seu tipo, emitindo poucos resíduos após sua queima. Ele pode ser encontrado em reservatórios no subsolo dos continentes e dos oceanos. Define-se o gás natural como a parte do petróleo que se encontra em forma de gás ou em solução, nas condições de temperatura e de pressão dos reservatórios, mas que em condições atmosféricas apresenta-se no estado gasoso. Quimicamente, o gás natural é uma mistura de hidrocarbonetos leves, sobretudo de metano (CH$_4$), etano (C$_2$H$_6$), propano (C$_3$H$_8$) e butano (C$_4$H$_{10}$). Mariana Lyra. Recursos Naturais: Gás Natural. \textbf{Revista Política Hoje}. 1ª Edição, V. 23, p. 149-174, 2014.

\par

No esquema abaixo, um dos componentes do gás natural participa de uma reação química. Átomos de diferentes elementos são representados por esferas de diferentes tamanhos:

\par

\begin{center}\(\underbrace{\mathrm{CH_4}}_{(1)} + \underbrace{2\,\mathrm{O_2}}_{(2)} \longrightarrow \underbrace{\mathrm{CO_2}}_{(3)} + \underbrace{2\,\mathrm{H_2O}}_{(4)}\)\end{center}

\par

Com base nas informações fornecidas e nos conhecimentos sobre ciências, assinale a afirmativa CORRETA:\begin{enumerate}
\item A reação representada é de combustão e (1) e (2) correspondem a substâncias simples.
\item (3) e (4) são representações de substâncias compostas e uma mistura gasosa de (3) e (4) será do tipo
      homogênea.
\item A reação é representada utilizando o modelo atômico de Dalton e (2) e (4) correspondem a substâncias
      compostas.
\item Os componentes do gás natural – metano, etano, propano e butano – podem ser separados por filtração.\end{enumerate}
